\section{O Método: Pipeline Núcleo-Transformação}
\label{sec:kernel_pipeline}

Esta seção descreve o método como um procedimento aplicado sobre uma amostra finita de $n$ objetos $z_i = (x_{i1},\dots,x_{id})$ do domínio produto $\mathcal{Z}$ da Equação \eqref{eq:product_domain}. Na aplicação a modelos de IA, esses objetos podem ser as entradas processadas pelo modelo, como as imagens utilizadas na Seção \ref{sec:metodologia}, ou unidades internas, como as posições espaciais de uma camada, mas a construção não depende dessa interpretação. Cada fonte $l$ é comparada por um \textit{kernel} positivamente definido $k_l$, e o \textit{kernel} produto $K$ da Equação \eqref{eq:tensor_sim_inner} realiza o tensor de similaridade da Definição \ref{def:tensor_sim}. Pelo Corolário \ref{cor:hadamard_form}, a sua matriz de Gram sobre a amostra, aqui denotada $\widetilde{\mathbf{G}}$ por ser anterior a qualquer transformação, é o produto de Hadamard
\begin{equation}\label{eq:raw_tensor_gram}
    \widetilde{\mathbf{G}} = \widetilde{\mathbf{G}}_1 \circ \cdots \circ \widetilde{\mathbf{G}}_d, \qquad (\widetilde{\mathbf{G}}_l)_{ij} = k_l(x_{il},x_{jl}),
\end{equation}
de modo que a construção tensorial permanece implícita e a interação entre as fontes é realizada pelo produto das Gram componentes.

\begin{definition}[Pipeline núcleo-transformação]
\label{def:kernel_pipeline}
Um \emph{pipeline núcleo-transformação} sobre $\mathcal{Z}$ é um par $\Pi = (K,T)$, em que $K$ é um \textit{kernel} positivamente definido sobre $\mathcal{Z}$ e $T\in\mathcal{T}_n$ é uma transformação pós-\textit{kernel} admissível (Definição \ref{def:T_admissible}). Sobre a amostra, $K$ produz a Gram bruta $\widetilde{G}_{ij} = K(z_i,z_j)$, e a matriz efetivamente utilizada nas leituras geométricas e espectrais é
\begin{equation}\label{eq:transformed_gram}
    \mathbf{G} = T(\widetilde{\mathbf{G}}).
\end{equation}
\end{definition}

O \textit{kernel} determina a geometria bruta do espaço de representação, e a transformação $T$ modifica a sua realização amostral, podendo remover componentes de translação, eliminar diferenças de escala, normalizar individualmente os objetos ou alterar a distribuição espectral da matriz. Diferentes escolhas de \textit{kernel} e de transformação não são, em geral, equivalentes: cada pipeline seleciona uma noção específica de estrutura geométrica. A matriz $\widetilde{\mathbf{G}}$ é a restrição amostral de $K$, enquanto $\mathbf{G}$ é, em geral, a Gram de uma geometria empírica transformada, que só é a restrição de um novo \textit{kernel} quando $T$ é induzida por uma transformação no nível do próprio \textit{kernel}. Essa hipótese não é necessária para as leituras desenvolvidas a seguir.

\begin{remark}[Transformações por componente]
\label{rem:component_transformations}
Quando diferentes propriedades devem ser impostas a cada fonte, podem ser utilizados pipelines componentes $\Pi_l = (k_l,T_l)$, com $\mathbf{G}_l = T_l(\widetilde{\mathbf{G}}_l)$ e matriz conjunta
\begin{equation}\label{eq:transformed_tensor_gram}
    \mathbf{G} = \mathbf{G}_1 \circ \cdots \circ \mathbf{G}_d,
\end{equation}
que permanece simétrica e semidefinida positiva pelo mesmo argumento do Corolário \ref{cor:hadamard_form}. Essa construção não coincide, em geral, com $T(\widetilde{\mathbf{G}}_1 \circ \cdots \circ \widetilde{\mathbf{G}}_d)$: a primeira impõe propriedades a cada fonte antes de combiná-las, e a segunda modifica diretamente a geometria conjunta.
\end{remark}

\subsubsection*{Axiomas para a transformação pós-\textit{kernel}}

\begin{definition}[Transformação pós-\textit{kernel} admissível]
\label{def:T_admissible}
Seja $\mathbb{S}_+^n$ o cone das matrizes simétricas semidefinidas positivas de ordem $n$. Uma aplicação $T: \mathbb{S}_+^n \to \mathbb{S}_+^n$ é dita uma \emph{transformação pós-\textit{kernel} admissível} se é equivariante a reindexações dos objetos, isto é, se
\begin{equation}\label{eq:T_permutation}
    T(P\widetilde{\mathbf{G}}P^\top) = P\,T(\widetilde{\mathbf{G}})\,P^\top
\end{equation}
para toda matriz de permutação $P$ e toda $\widetilde{\mathbf{G}}\in\mathbb{S}_+^n$. A classe dessas transformações é denotada por $\mathcal{T}_n$. Uma transformação admissível é dita \emph{normalizadora em traço} quando $\operatorname{tr}(T(\widetilde{\mathbf{G}})) = 1$ para toda matriz não nula do seu domínio de validade, e \emph{idempotente} quando $T\circ T = T$.
\end{definition}

\begin{remark}[Significado da admissibilidade]
\label{rem:T_axioms}
A condição $T: \mathbb{S}_+^n \to \mathbb{S}_+^n$ garante que $\mathbf{G}$ continue admitindo uma realização como produtos internos em algum espaço de características empírico. A equivariância \eqref{eq:T_permutation} garante que o resultado não dependa da ordem atribuída aos objetos, cuja numeração não tem, por si só, significado geométrico. A normalização e a idempotência não fazem parte da admissibilidade básica e apenas selecionam subclasses com comportamentos algébricos adicionais.
\end{remark}

A classe $\mathcal{T}_n$ é fechada por composição e contém as transformações espectrais
\begin{equation}\label{eq:spectral_transformation}
    T_f(\widetilde{\mathbf{G}}) = V f(\Lambda)V^\top, \qquad \widetilde{\mathbf{G}} = V\Lambda V^\top,
\end{equation}
com $f$ aplicada a cada autovalor e $f(\lambda)\geq 0$ para $\lambda\geq 0$, condição que preserva a semidefinição positiva. As leituras das seções seguintes dependem apenas de $\mathbf{G}\in\mathbb{S}_+^n$ e não exigem uma escolha particular de $T$, mas o seu conteúdo interpretativo é condicionado pelas propriedades que a transformação adotada preserva ou descarta.

\subsubsection*{Catálogo de transformações e propriedades geométricas}

Propriedades adicionais, como invariância a translações ou a reescalonamentos globais ou individuais, não decorrem da admissibilidade e devem ser verificadas para cada transformação. A Tabela \ref{tab:T_catalog} apresenta as transformações utilizadas neste trabalho, em que $H = I - \frac{1}{n}\mathbf{1}\mathbf{1}^{\top}$ é a matriz de centramento e $D = \operatorname{diag}(\widetilde{G}_{11},\dots,\widetilde{G}_{nn})$.

\begin{table}[pos=ht]
\scriptsize
\centering
\caption{Transformações pós-\textit{kernel}, suas invariâncias específicas e as estruturas geométricas enfatizadas. As transformações que envolvem divisões são definidas apenas quando os denominadores correspondentes são não nulos.}
\label{tab:T_catalog}
\renewcommand{\arraystretch}{1.4}
\begin{tabular}{p{1.5cm} p{2.5cm} p{4.3cm} p{4.6cm}}
\hline
\textbf{Símbolo}
&
\textbf{Definição}
&
\textbf{Transformação específica}
&
\textbf{Estrutura enfatizada}
\\
\hline

$T_c$
&
$H\widetilde{\mathbf{G}}H$
&
Translação comum das representações no espaço de características
&
Relações entre os objetos em torno do centroide
\\

$T_{\mathrm{tr}}$
&
$\displaystyle
\frac{\widetilde{\mathbf{G}}}
{\operatorname{tr}(\widetilde{\mathbf{G}})}$
&
Reescalonamento global positivo
&
Distribuição relativa da energia entre os modos espectrais
\\

$T_a$
&
$\displaystyle
D^{-1/2}
\widetilde{\mathbf{G}}
D^{-1/2}$
&
Reescalonamentos positivos individuais das representações
&
Similaridade angular, com descarte das magnitudes individuais
\\

$T_{\mathrm{tr}}\circ T_c$
&
$\displaystyle
\frac{
    H\widetilde{\mathbf{G}}H
}{
    \operatorname{tr}
    \left(
        H\widetilde{\mathbf{G}}H
    \right)
}$
&
Translação comum e reescalonamento global positivo
&
Geometria centrada expressa em escala relativa
\\

\hline
\end{tabular}
\end{table}

A transformação angular $T_a$ exige $\widetilde{G}_{ii}>0$ e tem entradas $\widetilde{G}_{ij}/\sqrt{\widetilde{G}_{ii}\widetilde{G}_{jj}}$, isto é, normaliza a representação de cada objeto e retém apenas a similaridade angular.

\begin{proposition}[Efeito de $T_c$ e $T_{\mathrm{tr}}$ sobre a geometria de distâncias]\label{prop:invariancia_distancias}
Seja $\mathbf{G}\in\mathbb{S}_+^n$ e sejam $d_{\mathbf{G}}^2(i,j)=G_{ii}+G_{jj}-2G_{ij}$ as distâncias ao quadrado induzidas (Seção \ref{sec:representation_geometry}). Então \textnormal{(i)} $d^2_{T_c(\mathbf{G})}(i,j)=d^2_{\mathbf{G}}(i,j)$ para todo $i,j$, e \textnormal{(ii)} $d^2_{T_{\mathrm{tr}}(\mathbf{G})}(i,j)=d^2_{\mathbf{G}}(i,j)/\operatorname{tr}(\mathbf{G})$. Em particular, as vizinhanças $N_k$, a ordenação dos vizinhos e qualquer leitura que dependa apenas da ordem das distâncias são invariantes a $T_c$ e a $T_{\mathrm{tr}}$ aplicadas à Gram conjunta.
\end{proposition}
\begin{proof}
Seja $\varphi_i$ uma representação de $\mathbf{G}$ no espaço de características, com $G_{ij}=\langle\varphi_i,\varphi_j\rangle$. A centragem $H\mathbf{G}H$ é a Gram de $\varphi_i-\bar\varphi$, com $\bar\varphi=\frac1n\sum_k\varphi_k$, e a diferença $\varphi_i-\varphi_j$ não se altera. A normalização por traço multiplica todas as entradas por $1/\operatorname{tr}(\mathbf{G})$.
\end{proof}
Segue que, nas leituras por vizinhança, margem e contrafactual, apenas as escolhas feitas \emph{antes} da composição (a transformação angular $T_a$ e o deslocamento por $\mathbf{J}$ em cada fonte, Subseção \ref{subsec:construcao}) alteram a geometria. As transformações $T_c$ e $T_{\mathrm{tr}}$, por sua vez, afetam o espectro e, portanto, as leituras de energia $\widehat{E}_{\mathrm{ratio}}$. A contribuição de $T_a$ e de $\mathbf{J}$ deve, assim, ser verificada por ablação (Subseção \ref{subsec:experimentos}), e não apenas assumida.

\subsubsection*{O que é fixo e o que é escolha}

A partir das definições apresentadas nesta seção, o método pode ser descrito por uma estrutura fixa e por cinco escolhas. A estrutura fixa corresponde à forma de composição, isto é, um \textit{kernel} positivamente definido para cada fonte, a combinação das fontes pelo produto de Hadamard e a aplicação de transformações pertencentes a $\mathcal{T}_n$. As escolhas, por sua vez, são: (i) a unidade indexada, isto é, o que cada objeto $z_i$ representa, (ii) as $d$ fontes de informação associadas a cada objeto, (iii) o \textit{kernel} $k_l$ de cada fonte, (iv) as transformações $T_l$ aplicadas a cada componente e a transformação $T$ aplicada à matriz conjunta, e (v) as leituras extraídas de $\mathbf{G}$, desenvolvidas na Seção \ref{sec:representation_geometry} e na Subseção \ref{subsec:spectral}. Veja que é a estrutura fixa que garante, para qualquer conjunto de escolhas, que $\mathbf{G}$ seja semidefinida positiva e, portanto, que ela admita a geometria e a decomposição espectral apresentadas nas seções seguintes. Essa garantia é de boa definição: ela não implica que a geometria resultante seja informativa sobre o comportamento do modelo, o que depende das fontes e dos \textit{kernels} escolhidos e é verificado empiricamente.

Essas escolhas não são neutras, visto que, como discutido anteriormente, cada uma delas seleciona uma noção específica de estrutura geométrica. Entretanto, elas não fazem parte do método. Um conjunto particular de escolhas é denominado \emph{instanciação}, e a Seção \ref{sec:metodologia} apresenta uma delas, composta por três fontes (ativações, gradientes contrastivos e parâmetros ativados) extraídas de redes neurais, e uma segunda, para árvores de decisão, com fontes próprias (caminho, folga e classes). Sendo assim, os resultados obtidos com essa instanciação avaliam as fontes e os \textit{kernels} escolhidos, enquanto aquilo que pode ser transferido para outros modelos e outras fontes é a forma do método.

O Algoritmo \ref{alg:pipeline} resume esse procedimento de forma genérica. Veja que as escolhas (i) a (iv) aparecem apenas como entradas do algoritmo, enquanto a estrutura fixa corresponde ao seu corpo.

\begin{algorithm}[htbp]
\caption{\textit{Pipeline} núcleo-transformação}
\label{alg:pipeline}
\begin{algorithmic}[1]
\Require objetos $z_i = (x_{i1},\dots,x_{id})$, $i = 1,\dots,n$, \textit{kernels} $k_1,\dots,k_d$, transformações $T_1,\dots,T_d, T \in \mathcal{T}_n$
\Ensure matriz de Gram $\mathbf{G} \in \mathbb{S}_+^n$
\State $\mathbf{G} \gets \mathbf{1}\mathbf{1}^\top$
\For{$l = 1,\dots,d$}
    \For{$i, j = 1,\dots,n$}
        \State $(\widetilde{\mathbf{G}}_l)_{ij} \gets k_l(x_{il}, x_{jl})$ \Comment{Gram bruta da fonte $l$}
    \EndFor
    \State $\mathbf{G} \gets \mathbf{G} \circ T_l(\widetilde{\mathbf{G}}_l)$ \Comment{produto de Hadamard}
\EndFor
\State \Return $T(\mathbf{G})$
\end{algorithmic}
\end{algorithm}